% Long Chain-of-Thought (LCoT) and Reinforcement Learning with Verifiable Rewards (RLVR) have been demonstrated to improve the reasoning capabilities of Large Language Models (LLMs). 
% We observe that graph-structured reasoning patterns emerging in LCoT, such as divide-and-conquer, aggregation, and reflection, contribute to the reasoning performance of LLMs. 
% To systematically exploit these patterns, we propose \textbf{GRP}, a \textbf{Graph Reasoning Paradigm} to model complex human cognition, and synthesize graph-structured data with step-level labels to train models, improving the interpretability, reliability, and accuracy of reasoning. 
% Furthermore, existing outcome-based reward methods lack fine-grained supervision and are susceptible to reward hacking, while process-based rewards suffer from costly data acquisition and limited verification effectiveness. 
% To address these issues, we propose \textbf{PASC-GRPO}, which enables \textbf{Process-Aware} evaluation through graph-structured outcome rewards without relying on process reward models and mitigates reward hacking via \textbf{Stratified Clipping Advantage Estimation}. Experiments demonstrate improvements in accuracy, interpretability, and training efficiency on mathematical reasoning and code generation tasks. Data, models, and code will be released later.

% Long Chain-of-Thought (LCoT) and Reinforcement Learning with Verifiable Rewards (RLVR) have proven to be effective in enhancing the reasoning capabilities of Large Language Models (LLMs). 
% However, existing reasoning steps are predominantly text-based, lacking traceability of reasoning paths and logical structure. 
% % Moreover, current outcome-based and process-based reward methods suffer from coarse-grained supervision, vulnerability to reward hacking, high data acquisition costs, and limited verification validity. 
% Moreover, outcome-based reward methods lack fine-grained supervision over the reasoning process and are vulnerable to reward hacking, while process-based reward methods suffer from substantial training overhead and limited generalization.
% To address these challenges, we propose a \textbf{G}raph \textbf{R}easoning \textbf{P}aradigm (\textbf{GRP}) to simulate complex human cognition by synthesizing graph-structured CoT with step-level annotations for supervised fine-tuning (SFT), thereby improving reasoning accuracy, interpretability, verifiability, and logical consistency. 
% We further present \textbf{P}rocess-\textbf{A}ware \textbf{S}tratified \textbf{C}lipping \textbf{GRPO} (\textbf{PASC-GRPO}), which achieves process-aware verification through graph-structured outcome rewards without relying on process reward models, and mitigates reward hacking via stratified clipping advantage estimation. 
% Experiments demonstrate significant improvements in accuracy across mathematical reasoning and code generation tasks. Data, models, and code will be released.

% Long Chain-of-Thought (LCoT) achieved by Reinforcement Learning with Verifiable Rewards (RLVR) have proven to be effective in enhancing the reasoning capabilities of Large Language Models (LLMs). However, 过去大模型的输出是plain-text，即使有RLVR从结果或者过程的层面进行优化，但是仍然面临着coarse-grained supervision，vulnerability to reward hacking，训练成本高，泛化性差的问题。如果推理过程是结构化、符号化的，就可以在计算结果奖励时低成本地把控推理过程。为此，我们提出了\textbf{G}raph \textbf{R}easoning \textbf{P}aradigm (\textbf{GRP})，尝试为复杂的推理过程标注细粒度标签，并建模成结构化的graph structure。并进一步设计了\textbf{P}rocess-\textbf{A}ware \textbf{S}tratified \textbf{C}lipping \textbf{GRPO} (\textbf{PASC-GRPO})，which achieves process-aware verification through graph-structured outcome rewards without relying on process reward models, and mitigates reward hacking via stratified clipping advantage estimation. Experiments demonstrate significant improvements in accuracy across mathematical reasoning and code generation tasks. Data, models, and code will be released.

Long Chain-of-Thought (LCoT), achieved by Reinforcement Learning with Verifiable Rewards (RLVR), has proven effective in enhancing the reasoning capabilities of Large Language Models (LLMs).
However, reasoning in current LLMs is primarily generated as plain text, where performing semantic evaluation on such unstructured data creates a computational bottleneck during training.
Despite RLVR-based optimization, existing methods still suffer from coarse-grained supervision, reward hacking, high training costs, and poor generalization.
To address these issues, we propose the Graph Reasoning Paradigm (\textbf{GRP}), which realizes structured and symbolic reasoning, implemented via graph-structured representations with step-level cognitive labels. 
Building upon GRP, we further design Process-Aware Stratified Clipping Group Relative Policy Optimization (\textbf{PASC-GRPO}), which leverages structured evaluation to replace semantic evaluation, achieves process-aware verification through graph-structured outcome rewards, and mitigates reward hacking via stratified clipping advantage estimation.
Experiments demonstrate significant improvements across mathematical reasoning and code generation tasks. Data, models, and code will be released later.

% However, existing LLMs output remain predominantly plain-text, and although with outcome- or process-level optimization via RLVR, current approaches still suffer from coarse-grained supervision, vulnerability to reward hacking, high training costs, and poor generalization.
% If the reasoning process is structured and symbolic, it becomes possible to exert fine-grained control over reasoning process at low cost when computing outcome rewards.